\chapter{Advanced Neural Networks}

\section{Convolutional Neural Networks (CNNs)}
Convolutional Neural Networks (CNNs) are a class of deep neural networks that have proven to be highly effective for tasks involving image and video recognition, as well as other applications that require spatial hierarchies of features. Before diving into the architecture of CNNs, it is essential to understand the fundamental components that make them unique.

\subsection{Convolution}

Convolution is the core operation in CNNs, which is used to extract features from input data. Humans can easily recognize objects in images by identifying patterns and features, such as edges, textures, and shapes. However, for a computer to recognize these features, it needs to learn how to detect them through a process called convolution.  It consists of:
\begin{itemize}
    \item The input image, which is typically represented as a multi-dimensional array (tensor).
    \item A filter (aka kernel), which is a smaller-sized matrix that is used to detect specific features in the input image.
    \item A \emph{feature map}, which is the output of the convolution operation, representing the presence of specific features in the input image.
\end{itemize}

In order to obtain the feature map, the filter is slid over the input image, and at each position, an \emph{element-wise} multiplication is performed between the filter and the corresponding region of the input image. The results are then summed up to produce a single value in the feature map. Some important parameters in the convolution operation include:
\begin{itemize}
    \item \ul{Stride}: The number of pixels by which the filter is moved across the input image. A larger stride results in a smaller feature map.
    \item \ul{Padding}: This process has two main problems: \begin{itemize}
        \item The size of the feature map is reduced after each convolution operation, which can lead to a loss of information, especially for deeper networks.
        \item The edges of the input image are not treated equally, as the filter has fewer neighboring pixels to convolve with at the borders, which can lead to artifacts in the feature map.
    \end{itemize}

    In order to address these issues, padding is applied to the input image before the convolution operation. It involves adding extra pixels (usually zeros) around the borders of the input image, allowing the filter to be applied to the edges and preserving the spatial dimensions of the feature map. Some common padding techniques include:
    \begin{itemize}
        \item \ul{Constant Padding}: Adding a constant value (usually zero) around the borders of the input image.
        \item \ul{Reflective Padding}: Mirroring the values of the input image at the borders.
        \item \ul{Replicate Padding}: Repeating the values of the input image at the borders.
    \end{itemize}
\end{itemize}

Some specific filters are:
\begin{itemize}
    \item \ul{Convolution Box Blur Filter}: This filter is used to blur an image by averaging the pixel values in a local neighborhood.
    \begin{equation*}
        \frac{1}{9}\begin{bmatrix}
            1 & 1 & 1 \\
            1 & 1 & 1 \\
            1 & 1 & 1
        \end{bmatrix}
    \end{equation*}

    \item \ul{Convolution Gaussian Blur Filter}: This filter is used to blur an image while preserving edges. It is based on the Gaussian function, which gives more weight to the central pixels in the local neighborhood.
    \begin{equation*}
        \frac{1}{16}\begin{bmatrix}
            1 & 2 & 1 \\
            2 & 4 & 2 \\
            1 & 2 & 1
        \end{bmatrix}
    \end{equation*}
    \item \ul{Convolution Edge Detection Filter}: This filter is used to detect edges in an image by highlighting areas of rapid intensity change.
    \begin{equation*}
        \begin{bmatrix}
            -1 & -1 & -1 \\
            -1 & 8 & -1 \\
            -1 & -1 & -1
        \end{bmatrix}
    \end{equation*}
\end{itemize}

\subsection{Pooling}

Pooling is a downsampling operation that reduces the spatial dimensions of the feature maps while retaining the most important information. It helps to reduce the computational complexity of the network and makes it more robust to variations in the input data, such as translations and distortions. There are several types of pooling operations, including:
\begin{itemize}
    \item \ul{Max Pooling}: This operation selects the maximum value from a local neighborhood of the feature map. It is the most commonly used pooling technique and helps to retain the most prominent features in the feature map.
    \item \ul{Average Pooling}: This operation computes the average value of a local neighborhood of the feature map. It is less commonly used than max pooling, as it tends to blur the features and may not preserve the most important information.
    \item \ul{ROI Pooling}: This operation is used in object detection tasks, where the input feature map is divided into a fixed number of Regions of Interest (ROIs), and \emph{each ROI is pooled to a fixed size}. This is useful for handling objects of varying sizes in the input image, as it allows the network to work with a fixed-size representation of the ROIs, regardless of their original size.
\end{itemize}

\subsection{Flattening}

Flattening is the process of converting the multi-dimensional feature maps into a one-dimensional vector, which can then be fed into fully connected layers for classification or regression tasks. This step is essential for transitioning from the convolutional and pooling layers to the dense layers of the network.

\subsection{Architecture of CNNs}

A typical CNN architecture is divided into several layers, each serving a specific purpose. The main layers in a CNN include:
\begin{itemize}
    \item \ul{Input Layer}: This layer receives the input data, which is usually an image represented as a multi-dimensional array (tensor). The input layer does not perform any computations but serves as the entry point for the data into the network.
    \item \ul{Convolutional Layers}: These layers perform the convolution operation, extracting features from the input data using \emph{learnable} filters.
    
    Each convolutional layer is typically followed by an activation function, such as ReLU (Rectified Linear Unit), which introduces non-linearity into the network and allows it to learn complex patterns in the data.

    \item \ul{Pooling Layers}: These layers perform the pooling operation, reducing the spatial dimensions of the feature maps and retaining the most important information.
    
    \item \ul{Flattening Layer}: After the convolutional and pooling layers, the feature maps are flattened into a one-dimensional vector, which can then be fed into the fully connected layers for classification or regression tasks.
    
    \item \ul{Fully Connected Layers}: Using the flattened feature vector, these layers perform the final classification or regression tasks. Each neuron in a fully connected layer is connected to every neuron in the previous layer, allowing the network to learn complex relationships between the features extracted by the convolutional and pooling layers.
    
    \item \ul{Output Layer}: This layer produces the final output of the network, which can be a class label for classification tasks or a continuous value for regression tasks. The output layer typically uses an activation function, such as softmax for multi-class classification or sigmoid for binary classification.
\end{itemize}


\section{Generative Adversarial Networks (GANs)}

GANs are a class of deep learning models that consist of two neural networks, a generator and a discriminator, which are trained simultaneously in a competitive setting.
\begin{itemize}
    \item The generator network takes random noise as input and generates synthetic data samples that resemble real images.
    \item The discriminator network takes both real and synthetic data samples as input and learns to distinguish between them, outputting a probability score indicating whether the input is real or fake.
\end{itemize}


\section{Autoencoders}

Autoencoders are a type of neural network used for unsupervised learning, where the goal is to reconstruct the input data as accurately as possible. They consist of:
\begin{itemize}
    \item An encoder network that compresses the input data into a lower-dimensional representation, called the latent space.
    \item The latent space representation, which captures the most important features of the input data while discarding noise and irrelevant information.
    \item A decoder network that reconstructs the original input data from the latent representation.
\end{itemize}

The concept of the latent space has a great potential, as latent variables can be combined to generate new data samples.